% One-page letter / statement template. Matches the CV letterhead.
\documentclass[11pt, letterpaper]{article}
\usepackage[T1]{fontenc}
\usepackage[utf8]{inputenc}
\usepackage{charter}
\usepackage[top=0.55in, bottom=0.45in, left=0.8in, right=0.8in]{geometry}
\usepackage[dvipsnames]{xcolor}
\definecolor{linkblue}{HTML}{1155CC}
\usepackage[colorlinks=true, urlcolor=linkblue, linkcolor=linkblue]{hyperref}
\usepackage{parskip}
\pagestyle{empty}
\setlength{\parindent}{0pt}
\linespread{1.0}

\newcommand{\letterhead}{%
  \begin{center}
    {\large\bfseries Rodrigo França Chaves}\\[2pt]
    {\small
      Chicago, IL \enspace$\cdot$\enspace
      \href{mailto:rchaves@uchicago.edu}{rchaves@uchicago.edu} \enspace$\cdot$\enspace
      +1 312 217 5526 \enspace$\cdot$\enspace
      \href{https://rodrigofch7.github.io}{rodrigofch7.github.io}
    }
  \end{center}
  \vspace{2pt}
  {\rule{\linewidth}{0.8pt}}
  \vspace{10pt}
}

\begin{document}
\letterhead

\textbf{Analysis Group} \\
Analyst, Generalist (2027 start) \\
\textit{Geographic preference: Washington, DC, first; Chicago, second.}

\vspace{8pt}

Dear Hiring Committee,

I am applying for the Analyst position starting in 2027. I am a Master's student in Computational Analysis and Public Policy at the University of Chicago, graduating June 2027, and I want to spend my career doing what Analysis Group does: taking a contested empirical question, building the evidence out of difficult data, and defending the answer to people who did not produce it. I interviewed with Analysis Group last year for a summer internship and reached the final round, during my first month in the United States; a year on, I am applying again with considerably more behind me.

I have already done a version of this work. I am first author of a peer-reviewed article in \textit{Utilities Policy} asking whether transferring Brazil's water and wastewater systems to private operators improved public health. It compares propensity-matched municipalities before and after transfer between 1998 and 2021 using staggered difference-in-differences, finding significant reductions in diarrheal morbidity among children under five and no effect on diseases unrelated to sanitation. I wrote the paper, assembled the data, and wrote the code. The result is genuinely mixed across municipalities, and saying so plainly was part of the work.

Much of my experience is in the conditions your cases involve: administrative records not designed for analysis. This summer at the World Bank I reconstructed four decades of Mumbai property tax policy into a single timeline, then tested millions of assessment records for bunching at policy thresholds, the pattern that appears when declared values are manipulated. Earlier, working with confidential São Paulo police microdata, the boundaries I needed were never publicly released, so I reconstructed them by clustering geocoded incidents into Voronoi polygons; without that step the analysis had no unit of observation.

I am also comfortable at scale and in code. I built a horizontal-equity audit of New York's property tax roll measuring each of 1.1 million properties against a peer group defined by borough, building class, tax class, decade built, size, and value band, rather than a citywide benchmark that would put a Manhattan co-op and a Queens two-family in the same comparison. Roughly 43 percent sit more than 15 percent from their peer-group median, though that definition and threshold were ours, so the figure measures dispersion within comparable groups rather than proven misassessment. I designed the methodology, built the pipeline, and trained the models. I work in Stata, Python, and R.

For two years at Insper I held research assistant and teaching assistant roles simultaneously, teaching econometrics and GIS to more than fifty students while running my own analyses. I also build the delivery end deliberately: both recent projects ship as dashboards, so someone who cannot read the code can still interrogate the result. Explaining a method to an audience with no reason to trust it yet is the part of this work I take most seriously.

The \textit{Utilities Policy} article was already published when we spoke in October, so it is not what has changed. Both the Mumbai work and the New York audit are new since then, as is a research assistantship at the Mansueto Institute building OCR pipelines over century-old archives.

Thank you for your consideration. I would welcome the chance to talk.

\vspace{6pt}

Sincerely, \\
Rodrigo França Chaves

\end{document}
